\chapter{Cómo hicimos los experimentos}
\label{cap:diseno}

\section{Instancias y experimentos}
Usamos grafos del FHCP Challenge Set. Son 1\,001 grafos, de 66 a 9\,528 vértices, y todos tienen
ciclo hamiltoniano \cite{aco9}. Esto es muy útil: como sabemos que la respuesta siempre es ``sí'',
podemos medir directamente qué tan seguido la encuentra cada algoritmo.

Hicimos tres experimentos (tabla~\ref{tab:disenos}). Una \emph{corrida} es la ejecución de un
algoritmo sobre un grafo con una semilla.

\begin{table}[H]
\centering\small\filas
\caption{Experimentos realizados.}
\label{tab:disenos}
\begin{tabularx}{\textwidth}{@{}lLLccL@{}}
\toprule
 & Objetivo & Grafos & Algoritmos & Semillas & Límite por corrida \\
\midrule
A & Comparar los seis algoritmos & graph4 a graph24 ($n=84$--$198$) & 8 & \semillasA & \tiempoA\,s o \evalsA{} evaluaciones \\
B & Benchmarking con el FHCP & graph1 a graph100 ($n=66$--$558$) & 2 & \semillasB & \tiempoB\,s o \evalsB{} evaluaciones \\
C & Probar grafos grandes & graph50, 200, 500 y 1001 ($n=348$--$9\,528$) & 8 & \semillasC & \tiempoC\,s o \evalsC{} evaluaciones \\
\bottomrule
\end{tabularx}
\end{table}

Contamos 8 algoritmos porque SA y ACO tienen dos versiones cada uno. En el experimento B solo
corrimos \alg{ACO + Pósa} y \alg{ABC}, que fueron los que mejor salieron en A. Así pudimos cubrir
100 grafos con un tiempo de cómputo razonable.

\section{Presupuesto de cada corrida}
Cada corrida termina en cuanto pasa una de estas tres cosas: encuentra un ciclo, se le acaban las
evaluaciones o se le acaba el tiempo. Al revisar los resultados nos dimos cuenta de que no todos
los algoritmos se detuvieron por la misma razón (tabla~\ref{tab:presupuesto}). SA usó sus \evalsA{}
evaluaciones en unos 35 a 40 segundos, porque cada paso es muy barato. En cambio, GA, ACO, DE y
PSO hacen más trabajo por evaluación y llegaron al límite de 60 segundos habiendo usado solo entre
400 mil y 1.2 millones de evaluaciones.

Por eso, la comparación del experimento A es \emph{a igual tiempo}, no a igual número de
evaluaciones. Nos parece la comparación más práctica, porque al final lo que importa es cuánto hay
que esperar, pero hay que tenerlo en cuenta al leer los resultados.

\begin{table}[H]
\centering\small\filas
\caption{Experimento A: evaluaciones usadas y por qué se detuvieron las corridas sin éxito.}
\label{tab:presupuesto}
\resizebox{\textwidth}{!}{% Generado por Reporte/actualizar.py; no editar a mano.
\begin{tabular}{lrrrrr}
\toprule
Algoritmo & Evaluaciones (mediana) & Corridas sin éxito & Paró por tiempo & Paró por evaluaciones & Tiempo promedio (s) \\
\midrule
GA & 1\,199\,256 & 210 & 210 & 0 & 60.0 \\
SA (swap) & 6\,000\,000 & 210 & 0 & 210 & 41.6 \\
SA (mixto) & 6\,000\,000 & 193 & 0 & 193 & 34.5 \\
ACO & 417\,665 & 210 & 210 & 0 & 60.0 \\
ACO + Pósa & 84\,170 & 151 & 151 & 0 & 45.3 \\
ABC & 4\,032\,394 & 148 & 128 & 20 & 44.1 \\
DE & 800\,358 & 210 & 210 & 0 & 60.0 \\
PSO & 572\,944 & 210 & 210 & 0 & 60.0 \\
\bottomrule
\end{tabular}
}
\end{table}

\paragraph{Equipo.} Una laptop con procesador \cpu{} (\hilos{} hilos), \sistema{} y Python
\versionpython{}. Para terminar en un tiempo razonable, corrimos varias corridas al mismo tiempo
(20 en el experimento A y 16 en el C). Este procesador combina núcleos rápidos y núcleos de
eficiencia, así que una corrida que cae en un núcleo lento alcanza a hacer menos evaluaciones en
el mismo tiempo. Esto mete algo de variación en los resultados, pero no afecta que sean correctos.

\section{Qué medimos}
\begin{itemize}
  \item \textbf{Tasa de éxito:} el porcentaje de corridas que encontraron un ciclo hamiltoniano.
        Como solo tenemos 10 semillas por grafo, la acompañamos de un intervalo de confianza del
        95\,\% (método de Wilson \cite{wilson1927}), que indica qué tanto podría variar ese
        porcentaje.
  \item \textbf{Grafos resueltos:} en cuántos grafos distintos hubo al menos una corrida exitosa.
  \item \textbf{Tiempo y evaluaciones} que tardó en encontrar el ciclo.
  \item \textbf{Aristas faltantes al final:} cuando no hay éxito, qué tan cerca se quedó.
  \item \textbf{Precisión, recall y F1.} El enunciado pide estas métricas, así que las adaptamos a
        nuestro caso. Un ``positivo'' es que el algoritmo diga que encontró un ciclo.
        \begin{itemize}
          \item La \emph{precisión} (de los ciclos que reporta, cuántos son reales) siempre es 1,
                porque cada ciclo se verifica antes de reportarlo y no puede haber falsos positivos.
          \item El \emph{recall} (de los grafos que sí tienen ciclo, en cuántos lo encontró) es
                igual a la tasa de éxito, porque todos los grafos del FHCP tienen ciclo.
          \item El \emph{F1} combina las dos y aquí depende solo del recall.
        \end{itemize}
        Por esto, en nuestro caso la métrica que realmente distingue a los algoritmos es la tasa
        de éxito.
\end{itemize}

\section{Cómo comprobamos que los resultados son correctos}
\label{sec:verificacion}
No quisimos confiar solo en lo que reporta cada algoritmo. Al terminar, revisamos todos los
resultados con un programa de verificación aparte, que no usa nada del código del proyecto: lee
los archivos de los grafos por su cuenta y vuelve a calcular todo a partir de las rutas guardadas.
La tabla~\ref{tab:verificacion} muestra lo que revisamos. Todo salió correcto. Las corridas que
no encontraron ciclo no son errores del programa: terminaron normalmente al acabarse su
presupuesto y devolvieron una ruta válida, solo que con aristas faltantes.

\begin{table}[H]
\centering\small\filas
\caption{Comprobaciones sobre las \the\numexpr\ejecA+\ejecB+\ejecC\relax{} corridas.}
\label{tab:verificacion}
\begin{tabularx}{\textwidth}{Lr}
\toprule
Comprobación & Resultado \\
\midrule
Están todas las corridas y ninguna repetida (A, B, C) & \ejecA, \ejecB, \ejecC \\
Cada ruta es una permutación válida de los vértices & todas \\
El costo guardado coincide con el que recalculamos & todas \\
Las corridas marcadas como éxito tienen costo 0, y solo esas & todas \\
Cada ciclo encontrado tiene su archivo \texttt{.tour} y es válido & \the\numexpr\exitosA+\exitosB+\exitosC\relax{} de \the\numexpr\exitosA+\exitosB+\exitosC\relax \\
Ninguna corrida se pasó de su límite de evaluaciones o de tiempo & todas \\
Las curvas de convergencia nunca empeoran y terminan en el costo final & todas \\
Prueba del verificador: acepta las soluciones oficiales del FHCP & 24 de 24 \\
Al repetir corridas exitosas con la misma semilla sale la misma ruta & sí \\
\bottomrule
\end{tabularx}
\end{table}
